% ============================================================================
% input.tex — Documento principal en sintaxis LaTeX2e
% Pipeline: input.tex -> input.tks (XML intermedio) -> input.svg
% ============================================================================

\documentclass[11pt,a4paper]{article}
\usepackage[utf8]{inputenc}
\margins{2.2cm}{2.2cm}{20mm}{20mm}

% --- Definición de macros y entornos personalizados en el preámbulo ---
\newcommand{\cuerpobase}{%
  \align{justify}%
  \setlength{\baselineskip}{5.5mm}%
  \setlength{\parindent}{8mm}%
  \parmargins{0mm}{0mm}%
  \parspacing{0mm}{2.4mm}%
  \rmfamily\mdseries\upshape\fontsize{11pt}\fontcolor{Black}%
}

\newcommand{\codigo}[1]{{\ttfamily\fontsize{9.5pt}\fontcolor{Teal}#1}}

\newcommand{\ficha}[2]{%
  \par{\raggedright\linespacing{5mm}\parmargins{6mm}{0mm}\parspacing{0.5mm}{1mm}%
  \sffamily\fontsize{10pt}\fontcolor{Navy}\textbf{#1:} %
  \rmfamily\fontsize{10.5pt}\fontcolor{Black}#2\par}%
}

\providecommand{\versionmotor}{0.2}
\def\etiquetarevision#1{[Rev.~#1]}

\pagestyle{headings}
\markboth{Manual de \LaTeXe{} en Teksto}{Arquitectura \TeX{} -> .tks -> .svg}
\makeindex

\newsavebox{\cajasello}
\sbox{\cajasello}{\fbox{\textbf{VALIDADO}}}

\newenvironment{aviso}[1][Aviso T\'ecnico]{%
  \par\raggedright\parmargins{8mm}{8mm}\parspacing{2mm}{2mm}%
  \sffamily\fontsize{10pt}\fontcolor{DarkSlateGray}\textbf{#1 --- }\itshape%
}{\par}

% --- Metadatos para \maketitle ---
\title{Demostraci\'on Integral de \LaTeXe{} en Teksto \etiquetarevision{\versionmotor}}
\author{Equipo de Desarrollo de Teksto \and Laboratorio de Tipograf\'{\i}a Digital}
\date{\today}
\thanks{Pipeline completo: entrada .tex -> representaci\'on intermedia .tks -> salida vectorial .svg.}

\begin{document}

\cuerpobase
\maketitle

\begin{abstract}
Este documento ejemplifica las funcionalidades implementadas en Teksto siguiendo el modelo de procesamiento de \TeX{} y \LaTeXe{}:
expansi\'on de macros (\codigo{\textbackslash newcommand}, \codigo{\textbackslash def}, \codigo{\textbackslash newenvironment}),
agrupamiento l\'exico con pila de estados, contadores jer\'arquicos, registros de longitud, selecci\'on de fuentes NFSS,
versalitas (\textsc{Small Caps}), super\'{\i}ndices y sub\'{\i}ndices, seccionado autom\'atico, listas anidadas,
entornos literales (\codigo{verbatim}), modo matem\'atico, gr\'aficos vectoriales (\codigo{picture}), referencias cruzadas y bibliograf\'{\i}a.
\end{abstract}

\tableofcontents

% Inclusión del capítulo externo capitulo1.tex
% ============================================================================
% capitulo1.tex — Módulo secundario incluido mediante % ============================================================================
% capitulo1.tex — Módulo secundario incluido mediante % ============================================================================
% capitulo1.tex — Módulo secundario incluido mediante \input{capitulo1}
% Ejemplifica el motor léxico TeX, símbolos, acentos, NFSS, tamaños y bloques.
% ============================================================================

\section{Motor L\'exico \TeX{}, S\'{\i}mbolos y Tipograf\'{\i}a}

Este fichero (\codigo{capitulo1.tex}) se incorpora din\'amicamente desde \codigo{input.tex} mediante \codigo{\textbackslash input\{capitulo1\}}.
En \LaTeXe{}, un salto de l\'{\i}nea simple equivale a un espacio, una l\'{\i}nea en blanco emite \codigo{\textbackslash par}, y los comentarios con \codigo{\%} descartan el resto de la l\'{\i}nea junto con los espacios iniciales de la siguiente%
---tal como define el motor de Knuth---evitando espacios indeseados.

\subsection{Acentos Componibles, Ligaduras y Caracteres Especiales}

El analizador soporta acentos con o sin llaves en cualquier vocal o consonante:
\'a, \`e, \^i, \"u, \~n, \=o, \.z, \u{g}, \v{s}, \H{o}, \c{c}, \k{a}, \r{a}, \b{b}, \d{a}, \t{oo},
adem\'as de letras europeas como \aa, \AA, \ae, \AE, \oe, \OE, \o, \O, \l, \L, \ss\ y sin punto (\i, \j).

Las ligaduras se transforman autom\'aticamente: ``comillas dobles inglesas'', `comillas simples',
intervalos num\'ericos con gui\'on medio (p\'ags.~15--30), incisos con raya larga---sin espacios---y signos de apertura ?`C\'omo opera el lexer? !`Con precisi\'on tipogr\'afica!

\ficha{Caracteres reservados}{\# \$ \% \& \_ \{ \} \~{} \^{} \textbackslash\ \textbar\ \textless\ \textgreater}
\ficha{S\'{\i}mbolos editoriales}{\dag\ \ddag\ \S\ \P\ \copyright\ \textregistered\ \texttrademark\ \pounds\ \euro\ \dots\ \textbullet}
\ficha{Logotipos y fecha}{\TeX, \LaTeX, \LaTeXe{} --- Compilado el \today}

\subsection{Selecci\'on de Fuentes NFSS y Escala de Tama\~nos}

\ficha{Familias y series}{\textrm{Roman Serif}, \textsf{Sans Serif}, \texttt{Typewriter Mono}, \textbf{Negrita}, \textit{Cursiva}, \textsc{Versalitas}, \underline{Subrayado}}
\ficha{Super y sub\'{\i}ndices}{F\'ormula qu\'{\i}mica H\textsubscript{2}SO\textsubscript{4}, energ\'{\i}a E=mc\textsuperscript{2}, secci\'on 1.\textsuperscript{a}}
\ficha{Primitivas NFSS}{{\usefont{OT1}{cmss}{bx}{it} Fuente activada con \textbackslash usefont\{OT1\}\{cmss\}\{bx\}\{it\}}}
\ficha{Escala est\'andar}{{\tiny tiny} {\scriptsize scriptsize} {\footnotesize footnotesize} {\small small} {\normalsize normalsize} {\large large} {\Large Large} {\LARGE LARGE} {\huge huge}}

\subsubsection{Subsecci\'on de Tercer Nivel y T\'{\i}tulos en L\'{\i}nea}

\paragraph{P\'arrafo con t\'{\i}tulo (\textbackslash paragraph)} El comando \codigo{\textbackslash paragraph} compone su t\'{\i}tulo en negrita en la misma l\'{\i}nea de inicio del p\'arrafo, seguido de un espacio horizontal \codigo{\textbackslash quad}.

\begin{aviso}[Nota de Arquitectura]
El fichero intermedio generado por la primera etapa lleva ahora la extensi\'on \codigo{.tks} (formato estructurado XML de Teksto), a partir del cual \codigo{xml2svg} produce el documento interactivo \codigo{.svg}.
\end{aviso}

% Ejemplifica el motor léxico TeX, símbolos, acentos, NFSS, tamaños y bloques.
% ============================================================================

\section{Motor L\'exico \TeX{}, S\'{\i}mbolos y Tipograf\'{\i}a}

Este fichero (\codigo{capitulo1.tex}) se incorpora din\'amicamente desde \codigo{input.tex} mediante \codigo{\textbackslash input\{capitulo1\}}.
En \LaTeXe{}, un salto de l\'{\i}nea simple equivale a un espacio, una l\'{\i}nea en blanco emite \codigo{\textbackslash par}, y los comentarios con \codigo{\%} descartan el resto de la l\'{\i}nea junto con los espacios iniciales de la siguiente%
---tal como define el motor de Knuth---evitando espacios indeseados.

\subsection{Acentos Componibles, Ligaduras y Caracteres Especiales}

El analizador soporta acentos con o sin llaves en cualquier vocal o consonante:
\'a, \`e, \^i, \"u, \~n, \=o, \.z, \u{g}, \v{s}, \H{o}, \c{c}, \k{a}, \r{a}, \b{b}, \d{a}, \t{oo},
adem\'as de letras europeas como \aa, \AA, \ae, \AE, \oe, \OE, \o, \O, \l, \L, \ss\ y sin punto (\i, \j).

Las ligaduras se transforman autom\'aticamente: ``comillas dobles inglesas'', `comillas simples',
intervalos num\'ericos con gui\'on medio (p\'ags.~15--30), incisos con raya larga---sin espacios---y signos de apertura ?`C\'omo opera el lexer? !`Con precisi\'on tipogr\'afica!

\ficha{Caracteres reservados}{\# \$ \% \& \_ \{ \} \~{} \^{} \textbackslash\ \textbar\ \textless\ \textgreater}
\ficha{S\'{\i}mbolos editoriales}{\dag\ \ddag\ \S\ \P\ \copyright\ \textregistered\ \texttrademark\ \pounds\ \euro\ \dots\ \textbullet}
\ficha{Logotipos y fecha}{\TeX, \LaTeX, \LaTeXe{} --- Compilado el \today}

\subsection{Selecci\'on de Fuentes NFSS y Escala de Tama\~nos}

\ficha{Familias y series}{\textrm{Roman Serif}, \textsf{Sans Serif}, \texttt{Typewriter Mono}, \textbf{Negrita}, \textit{Cursiva}, \textsc{Versalitas}, \underline{Subrayado}}
\ficha{Super y sub\'{\i}ndices}{F\'ormula qu\'{\i}mica H\textsubscript{2}SO\textsubscript{4}, energ\'{\i}a E=mc\textsuperscript{2}, secci\'on 1.\textsuperscript{a}}
\ficha{Primitivas NFSS}{{\usefont{OT1}{cmss}{bx}{it} Fuente activada con \textbackslash usefont\{OT1\}\{cmss\}\{bx\}\{it\}}}
\ficha{Escala est\'andar}{{\tiny tiny} {\scriptsize scriptsize} {\footnotesize footnotesize} {\small small} {\normalsize normalsize} {\large large} {\Large Large} {\LARGE LARGE} {\huge huge}}

\subsubsection{Subsecci\'on de Tercer Nivel y T\'{\i}tulos en L\'{\i}nea}

\paragraph{P\'arrafo con t\'{\i}tulo (\textbackslash paragraph)} El comando \codigo{\textbackslash paragraph} compone su t\'{\i}tulo en negrita en la misma l\'{\i}nea de inicio del p\'arrafo, seguido de un espacio horizontal \codigo{\textbackslash quad}.

\begin{aviso}[Nota de Arquitectura]
El fichero intermedio generado por la primera etapa lleva ahora la extensi\'on \codigo{.tks} (formato estructurado XML de Teksto), a partir del cual \codigo{xml2svg} produce el documento interactivo \codigo{.svg}.
\end{aviso}

% Ejemplifica el motor léxico TeX, símbolos, acentos, NFSS, tamaños y bloques.
% ============================================================================

\section{Motor L\'exico \TeX{}, S\'{\i}mbolos y Tipograf\'{\i}a}

Este fichero (\codigo{capitulo1.tex}) se incorpora din\'amicamente desde \codigo{input.tex} mediante \codigo{\textbackslash input\{capitulo1\}}.
En \LaTeXe{}, un salto de l\'{\i}nea simple equivale a un espacio, una l\'{\i}nea en blanco emite \codigo{\textbackslash par}, y los comentarios con \codigo{\%} descartan el resto de la l\'{\i}nea junto con los espacios iniciales de la siguiente%
---tal como define el motor de Knuth---evitando espacios indeseados.

\subsection{Acentos Componibles, Ligaduras y Caracteres Especiales}

El analizador soporta acentos con o sin llaves en cualquier vocal o consonante:
\'a, \`e, \^i, \"u, \~n, \=o, \.z, \u{g}, \v{s}, \H{o}, \c{c}, \k{a}, \r{a}, \b{b}, \d{a}, \t{oo},
adem\'as de letras europeas como \aa, \AA, \ae, \AE, \oe, \OE, \o, \O, \l, \L, \ss\ y sin punto (\i, \j).

Las ligaduras se transforman autom\'aticamente: ``comillas dobles inglesas'', `comillas simples',
intervalos num\'ericos con gui\'on medio (p\'ags.~15--30), incisos con raya larga---sin espacios---y signos de apertura ?`C\'omo opera el lexer? !`Con precisi\'on tipogr\'afica!

\ficha{Caracteres reservados}{\# \$ \% \& \_ \{ \} \~{} \^{} \textbackslash\ \textbar\ \textless\ \textgreater}
\ficha{S\'{\i}mbolos editoriales}{\dag\ \ddag\ \S\ \P\ \copyright\ \textregistered\ \texttrademark\ \pounds\ \euro\ \dots\ \textbullet}
\ficha{Logotipos y fecha}{\TeX, \LaTeX, \LaTeXe{} --- Compilado el \today}

\subsection{Selecci\'on de Fuentes NFSS y Escala de Tama\~nos}

\ficha{Familias y series}{\textrm{Roman Serif}, \textsf{Sans Serif}, \texttt{Typewriter Mono}, \textbf{Negrita}, \textit{Cursiva}, \textsc{Versalitas}, \underline{Subrayado}}
\ficha{Super y sub\'{\i}ndices}{F\'ormula qu\'{\i}mica H\textsubscript{2}SO\textsubscript{4}, energ\'{\i}a E=mc\textsuperscript{2}, secci\'on 1.\textsuperscript{a}}
\ficha{Primitivas NFSS}{{\usefont{OT1}{cmss}{bx}{it} Fuente activada con \textbackslash usefont\{OT1\}\{cmss\}\{bx\}\{it\}}}
\ficha{Escala est\'andar}{{\tiny tiny} {\scriptsize scriptsize} {\footnotesize footnotesize} {\small small} {\normalsize normalsize} {\large large} {\Large Large} {\LARGE LARGE} {\huge huge}}

\subsubsection{Subsecci\'on de Tercer Nivel y T\'{\i}tulos en L\'{\i}nea}

\paragraph{P\'arrafo con t\'{\i}tulo (\textbackslash paragraph)} El comando \codigo{\textbackslash paragraph} compone su t\'{\i}tulo en negrita en la misma l\'{\i}nea de inicio del p\'arrafo, seguido de un espacio horizontal \codigo{\textbackslash quad}.

\begin{aviso}[Nota de Arquitectura]
El fichero intermedio generado por la primera etapa lleva ahora la extensi\'on \codigo{.tks} (formato estructurado XML de Teksto), a partir del cual \codigo{xml2svg} produce el documento interactivo \codigo{.svg}.
\end{aviso}


\section{Listas Anidadas, Texto Literal y Espaciado}

En esta segunda p\'agina (p\'agina actual: \thepage), maquetada a dos columnas mediante \codigo{\textbackslash twocolumn}, se verifican los entornos de listas de hasta cuatro niveles, el modo literal y los rellenos el\'asticos.

\subsection{Entornos de Listas}

\begin{enumerate}
  \item Primer elemento numerado del nivel principal (\codigo{enumi}).
  \item Segundo elemento con sublista anidada:
  \begin{enumerate}
    \item Sub-elemento alfab\'etico de nivel 2 (\codigo{enumii}).
    \item Sub-elemento con lista de vi\~netas interna:
    \begin{itemize}
      \item Vi\~neta est\'andar definida por \codigo{\textbackslash labelitemi}.
      \item[\ddag] Vi\~neta personalizada con argumento opcional \codigo{\textbackslash item[\textbackslash ddag]}.
    \end{itemize}
  \end{enumerate}
  \item Tercer elemento del nivel principal.
\end{enumerate}

\begin{description}
  \item[Entrada (.tex):] C\'odigo fuente en sintaxis est\'andar \LaTeXe{}.
  \item[Intermedio (.tks):] Representaci\'on estructurada XML de estilos, marcos y flujo.
  \item[Salida (.svg):] Documento vectorial multip\'agina interactivo.
\end{description}

\subsection{Modo Literal (\textbackslash verb y verbatim)}

El comando en l\'{\i}nea \verb|\newcommand{\cmd}[2][def]{#1-#2}| conserva intactos los caracteres especiales, mientras que \verb*|\verb* muestra espacios| visibiliza cada espacio con \verb*| |.

\begin{verbatim}
% Bloque literal verbatim:
\def\identidad#1{#1}
  Sangría y   espacios   preservados.
\end{verbatim}

\subsection{Saltos, Rellenos El\'asticos y Reglas}

Salto de l\'{\i}nea manual con \codigo{\textbackslash\textbackslash}:\\
Segunda l\'{\i}nea sin justificaci\'on forzada.\\[2.5mm]
Tercera l\'{\i}nea tras un espacio adicional de \codigo{2.5mm}.

\medskip
\noindent Ejemplos de rellenos horizontales y reglas (\codigo{\textbackslash hfill}, \codigo{\textbackslash dotfill}, \codigo{\textbackslash hrulefill}, \codigo{\textbackslash rule}):\\
Extremo izquierdo\hfill Extremo derecho\\
Secci\'on 1\dotfill P\'ag.~1\\
Secci\'on 2\dotfill P\'ag.~2\\
Firma autorizada:\hrulefill\ \rule{10mm}{1.8mm}

\bigskip
\section{Contadores, Longitudes y Ap\'endices}

\newcounter{experimento}[section]
\stepcounter{experimento}
\addtocounter{experimento}{2}

El contador de usuario \codigo{experimento} (dependiente de \codigo{section}) vale ahora \arabic{experimento} en ar\'abigo, \roman{experimento} en romano min\'uscula, \Roman{experimento} en romano may\'uscula y \Alph{experimento} en alfab\'etico.

\subsection{Aritm\'etica de Longitudes y Grupos}

\newlength{\margenlocal}
\setlength{\margenlocal}{1.2em}
\addtolength{\margenlocal}{4pt}

{\setlength{\parindent}{\margenlocal}
Este p\'arrafo modifica localmente \codigo{\textbackslash parindent} sumando unidades relativas (\codigo{1.2em}) y absolutas (\codigo{4pt}). Al cerrarse el grupo de llaves, el valor original de \codigo{\textbackslash parindent} se restaura autom\'aticamente.\par}

\section*{Secci\'on sin Numerar (\textbackslash section*)}

La variante con estrella \codigo{\textbackslash section*} aplica el estilo tipogr\'afico de secci\'on sin incrementar ni imprimir el contador.

\section{Cajas LR, Tablas y Elementos Flotantes}
\index{cajas LR}\index{tablas tabular}

Ejemplos de cajas horizontales: \mbox{ irrompible }, \fbox{con borde}, \framebox[24mm][c]{\textbf{centrada}}, desplazamiento vertical con \raisebox{1.2mm}{\small arriba} y \raisebox{-1.0mm}{\small abajo}, y caja prealmacenada: \usebox{\cajasello}.

\begin{table}[htb]
\begin{center}
\begin{tabular}{||l|c|r||}
\hline\hline
\multicolumn{3}{||c||}{\textbf{Resumen del Pipeline de Teksto}} \\
\hline
\textbf{Etapa} & \textbf{Formato} & \textbf{Salida} \\
\cline{2-3}
Frontend \TeX{} & \codigo{.tex} & \codigo{input.tks} \\
Motor Gr\'afico & \codigo{.tks} & \codigo{input.svg} \\
\hline\hline
\end{tabular}
\end{center}
\caption{Tabla con bordes dobles, \textbackslash multicolumn y \textbackslash cline.}\label{tab:pipeline}
\end{table}

\section{Modo Matem\'atico, Gr\'aficos Vectoriales y Referencias}\label{sec:matematicas}
\index{modo matematico}\index{entorno picture}

Como se detall\'o en la Secci\'on~\ref{sec:lexico} (p\'agina~\pageref{sec:lexico}) y se resumi\'o en el Cuadro~\ref{tab:pipeline}, el motor procesa el documento en dos pasadas para resolver referencias cruzadas y citas bibliogr\'aficas como \cite{knuth84,lamport94}.

En modo matem\'atico en l\'{\i}nea se componen expresiones como $E = mc^2$, desigualdades $\alpha + \beta \geq \gamma$, vectores $\vec{v} \cdot \overline{w} = 0$ y l\'{\i}mites $\lim_{n \to \infty} (1 + 1/n)^n = e$. A continuaci\'on se muestra la ecuaci\'on numerada~(\ref{eq:gauss}) con fracciones, radicales, coeficientes binomiales y operadores:

\begin{equation}\label{eq:gauss}
\int_{-\infty}^{\infty} e^{-x^2}\,dx = \sqrt{\pi}, \qquad \binom{n}{k} = \frac{n!}{k!(n-k)!}, \qquad \widehat{f}(\xi) = \sum_{k=0}^{N} \frac{\sqrt[3]{a_k + 1}}{2^k}
\end{equation}

\subsection{Diagrama Vectorial con el Entorno \codigo{picture}}

\begin{center}
\setlength{\unitlength}{1mm}
\begin{picture}(110,28)(0,0)
  \thinlines
  \put(4,14){\framebox(24,10)[c]{\textbf{.tex}}}
  \put(28,19){\vector(1,0){14}}
  \put(42,14){\dashbox{1.5}(26,10)[c]{\textbf{.tks (XML)}}}
  \put(68,19){\vector(1,0){14}}
  \put(82,14){\framebox(24,10)[c]{\textbf{.svg}}}
  \put(55,6){\oval(96,8)}
  \put(16,4){\makebox(78,4)[c]{\footnotesize Pipeline vectorial nativo de Teksto}}
  \put(10,6){\circle*{2.5}}
  \put(100,6){\circle{3.5}}
  \qbezier(16,14)(55,0)(94,14)
\end{picture}
\end{center}

\appendix
\section{Especificaci\'on del Fichero Intermedio \codigo{.tks}}

Tras ejecutar \codigo{\textbackslash appendix}, las secciones principales se numeran alfab\'eticamente (\thesection). El fichero \codigo{input.tks} contiene la representaci\'on XML intermedia con las definiciones de p\'agina, familias tipogr\'aficas, estilos deduplicados (\codigo{<para-style>}, \codigo{<char-style>}) y primitivas de flujo (\codigo{<p>}, \codigo{<span>}, \codigo{<br/>}, \codigo{<hspace/>}, \codigo{<vspace/>}, \codigo{<rule/>}, \codigo{<mfrac/>}, \codigo{<msqrt/>}, \codigo{<mover/>}, \codigo{<table>}, \codigo{<picture>}).

\begin{thebibliography}{9}
  \bibitem{knuth84} Donald E. Knuth, \textit{The \TeX book}, Addison-Wesley, 1984.
  \bibitem{lamport94} Leslie Lamport, \textit{\LaTeX: A Document Preparation System}, 2.\textsuperscript{a} ed., Addison-Wesley, 1994.
\end{thebibliography}

\printindex

\end{document}

